\chapter{Trichomoniasis}
\index{Trichomoniasis}

\section*{Definition and Overview}
Trichomoniasis (often referred to colloquially as ``trich'') is a highly prevalent, non-viral sexually transmitted infection (STI) caused by the flagellated protozoan parasite \textit{Trichomonas vaginalis}. Although often perceived as a minor or transient infection due to the high proportion of asymptomatic cases, it represents a substantial global public health burden and is classified by the World Health Organization (WHO) as the most common curable non-viral STI worldwide. The primary site of infection in females is the lower female genital tract, including the vulva, vagina, cervix, and urethra, while in males, it predominantly colonizes the urethra, prostate gland, and epididymis.

The clinical and epidemiological significance of trichomoniasis extends far beyond localized urogenital discomfort. The infection causes robust mucosal inflammation, which disrupts the epithelial barrier and recruits key immune cells, thereby significantly increasing the risk of acquiring and transmitting human immunodeficiency virus (HIV) and other sexually transmitted infections. Furthermore, maternal infection during pregnancy is associated with adverse obstetric outcomes, including preterm birth, low birth weight, and premature rupture of membranes. Despite its high prevalence and associated clinical risks, trichomoniasis remains relatively neglected in comparison to other STIs like chlamydia and gonorrhea. This neglect is partly due to the lack of routine screening programs, widespread reliance on low-sensitivity diagnostic methods like wet mount microscopy, and the historically false assumption that the infection is largely benign. This chapter provides a comprehensive, authoritative, and educational reference on the epidemiology, pathophysiology, clinical presentation, diagnosis, and management of trichomoniasis.

\section*{Epidemiology}
The epidemiological landscape of trichomoniasis is characterized by high global prevalence, striking demographic variations, and significant disparities in healthcare access. The World Health Organization estimates that there are over 156 million new cases of trichomoniasis worldwide each year, representing nearly half of all curable STIs globally. Unlike chlamydia and gonorrhea, which primarily affect adolescents and young adults under the age of 25, the prevalence of trichomoniasis tends to remain stable or even increase in older populations, particularly women over the age of 30 and 40. This unique age distribution suggests that the infection can persist as a chronic, low-grade condition for years if left untreated, forming a persistent reservoir within communities.

Demographic studies reveal a higher rate of diagnosed infections in women than in men. This difference is partly due to anatomical variations and clinical presentation; men are more likely to have transient, self-limiting infections or asymptomatic urethral colonization that goes undetected. High-risk populations include individuals residing in socioeconomically disadvantaged areas, those with multiple concurrent or consecutive sexual partners, individuals with a history of other STIs, and incarcerated populations. In many developed nations, substantial racial and ethnic disparities exist, with certain minority groups experiencing a disproportionately high burden of infection, often linked to systemic barriers to healthcare, unequal access to quality sexual health services, and higher background rates of the infection within specific sexual networks.

Geographically, the burden of trichomoniasis is heaviest in resource-limited regions, including sub-Saharan Africa, parts of Latin America, and South-East Asia. In these areas, diagnostic infrastructure is often lacking, and syndromic management protocols may fail to differentiate trichomoniasis from other causes of vaginal discharge or urethritis, leading to under-diagnosis and inadequate treatment. The primary mode of transmission is almost exclusively through direct sexual contact, including vaginal-penile intercourse and vaginal-vaginal contact. While non-sexual transmission via fomites (such as shared wet towels, toilet seats, or bathwater) has been historically debated due to the parasite's ability to survive in moist environments for short periods, it is clinically negligible and not considered a primary vector of transmission.

\section*{Aetiology and Pathophysiology}
The causative agent of trichomoniasis is \textit{Trichomonas vaginalis}, an obligate, extracellular, flagellated protozoan parasite. Structurally, the parasite exists solely in the active, motile trophozoite stage; it does not form a protective cyst stage, which is why it relies on direct, close physical contact for transmission. The trophozoite is typically pear-shaped (pyriform) but can assume an amoeboid shape when adhering to epithelial cells. It possesses five flagella: four anterior flagella that emerge from a single canal and a fifth flagellum that runs posteriorly along the outer edge of an undulating membrane. A rigid microtubular structure called the axostyle runs through the center of the organism and projects from the posterior end, likely assisting in attachment to host cells.

The pathophysiology of \textit{Trichomonas vaginalis} infection involves a complex sequence of cellular interactions, mechanical damage, and immunological activation:
\begin{itemize}
    \item \textbf{Adhesion and Colonization:} Upon entering the urogenital tract, the parasite must adhere to the squamous epithelial cells of the vaginal or urethral mucosa. This attachment is mediated by specific surface proteins called adhesins (such as AP65, AP51, AP33, and AP23) and lipophosphoglycans (LPG) present on the protozoan's cell surface. This binding triggers a morphological transition from the flagellated pyriform shape to a flattened, amoeboid shape, maximizing the surface contact area with the host cell.
    \item \textbf{Cytotoxicity and Tissue Damage:} Once adhered, \textit{T. vaginalis} exerts direct cytotoxic effects on the host epithelium. It secretes hemolytic toxins, pore-forming proteins, and a variety of cysteine peptidases (proteolytic enzymes) that degrade host extracellular matrix proteins, including collagen, fibronectin, and laminin. The parasite also phagocytoses host cells, including epithelial cells, red blood cells, and the protective lactobacilli of the normal vaginal flora.
    \item \textbf{Microbial Dysbiosis:} By consuming lactobacilli, the parasite disrupts the normal vaginal microbiome. Lactobacilli are responsible for maintaining an acidic vaginal pH (3.8 to 4.5) through the production of lactic acid. The depletion of these bacteria leads to an elevation of vaginal pH (frequently greater than 4.5), creating an environment that favors the survival of the parasite and promotes the growth of opportunistic anaerobic bacteria, often leading to co-existing bacterial vaginosis.
    \item \textbf{Inflammatory Response:} The host immune system detects the parasite via toll-like receptors (TLRs), triggering the release of pro-inflammatory cytokines and chemokines (such as IL-8 and TNF-alpha) from epithelial cells. This initiates a robust cellular immune response characterized by the infiltration of polymorphonuclear neutrophils, macrophages, and T-lymphocytes to the site of infection. This intense localized inflammation causes tissue edema, mucosal erythema, and micro-hemorrhages, which are clinically visible as diffuse red spots on the cervix (the ``strawberry cervix'').
    \item \textbf{Risk Factors for Infection:} Key risk factors that increase susceptibility or facilitate transmission include:
    \begin{itemize}
        \item \textbf{Unprotected Sexual Intercourse:} Engaging in sexual activity without barrier methods (such as latex condoms) allows direct mucosal exposure to the parasite.
        \item \textbf{Multiple or High-Risk Sexual Partners:} Having multiple partners increases the statistical likelihood of encountering an infected individual.
        \item \textbf{History of Prior STIs:} Previous infections may compromise mucosal integrity or indicate high-risk sexual networks.
        \item \textbf{Vaginal Douching:} The use of intravaginal cleansers alters the natural vaginal flora and elevates pH, compromising the biological barriers against colonization.
        \item \textbf{Socioeconomic Deprivation:} Limited access to healthcare services, sexual health education, and affordable diagnostic testing delays diagnosis and treatment, promoting transmission.
    \end{itemize}
\end{itemize}

\section*{Clinical Features}
The clinical presentation of trichomoniasis is highly variable, ranging from complete absence of symptoms to severe, acute inflammatory vaginitis or urethritis. The incubation period typically spans from 5 to 28 days post-exposure, though many individuals remain asymptomatic for months or even years.

In females, approximately 70\% to 85\% of infections are asymptomatic or minimally symptomatic, acting as a silent reservoir for transmission. When symptoms do manifest, they can be highly distressing and include:
\begin{itemize}
    \item \textbf{Vaginal Discharge:} The classic discharge associated with trichomoniasis is profuse, thin, purulent, and frothy or bubbly. It is typically yellow-green or greenish-yellow in color and is often accompanied by a strong, unpleasant, musty, or ``fishy'' odor.
    \item \textbf{Vulvovaginal Irritation:} Patients frequently experience intense itching (pruritus), burning, soreness, and redness of the vulva, perineum, and inner thighs.
    \item \textbf{Dysuria and Urinary Symptoms:} If the parasite colonizes the urethra or bladder, it can cause burning during urination, increased urinary frequency, and urgency.
    \item \textbf{Dyspareunia:} Superficial pain or discomfort during sexual intercourse is common due to severe mucosal inflammation.
    \item \textbf{Pelvic Discomfort:} Mild lower abdominal pain or pelvic heaviness may be reported.
    \item \textbf{Physical Examination Findings:} A speculum examination typically reveals erythema of the vaginal walls and cervix. In approximately 2\% to 5\% of cases (and up to 45\% when viewed under colposcopic magnification), punctate, macular mucosal hemorrhages are visible on the cervix, presenting a classic ``strawberry cervix'' (\textit{colpitis macularis}) appearance.
\end{itemize}

In males, the vast majority of infections (exceeding 70\% to 80\%) are entirely asymptomatic. When symptoms occur, they are usually transient and self-limiting, which contributes to the under-diagnosis and lack of treatment in men. Symptoms in males include:
\begin{itemize}
    \item \textbf{Urethral Discharge:} A mild, thin, mucoid, or purulent discharge from the penis, most noticeable in the morning.
    \item \textbf{Urethritis and Dysuria:} A burning sensation during urination or a feeling of itching and irritation inside the urethra.
    \item \textbf{Post-Ejaculatory Discomfort:} A burning or painful sensation immediately following ejaculation.
    \item \textbf{Complications:} In rare instances, the parasite can ascend to the prostate or epididymis, causing symptoms of prostatitis (perineal pain, dysuria, pelvic discomfort) or epididymitis (unilateral testicular pain and swelling).
\end{itemize}

\section*{Differential Diagnosis}
Given that the symptoms of trichomoniasis (such as abnormal discharge, itching, and dysuria) overlap significantly with other common urogenital conditions, clinical diagnosis based on symptoms alone is highly inaccurate. It is essential to differentiate trichomoniasis from other causes of vaginitis, cervicitis, and urethritis:
\begin{itemize}
    \item \textbf{Bacterial Vaginosis (BV):} An dysbiosis characterized by a reduction in hydrogen peroxide-producing lactobacilli and an overgrowth of anaerobic bacteria (e.g., \textit{Gardnerella vaginalis}, \textit{Atopobium vaginae}). Like trichomoniasis, BV is associated with an elevated vaginal pH (greater than 4.5) and a foul-smelling vaginal discharge. However, BV discharge is typically thin, homogenous, and grayish-white, rather than yellow-green and frothy. Furthermore, BV characteristically lacks signs of vulvovaginal inflammation (erythema, itching, dyspareunia), and microscopic examination reveals ``clue cells'' (epithelial cells coated with bacteria) without motile protozoa.
    \item \textbf{Vulvovaginal Candidiasis (Yeast Infection):} A fungal infection caused by \textit{Candida} species, most commonly \textit{Candida albicans}. It presents with intense vulvar itching, burning, and erythema, which can mimic trichomoniasis. However, the discharge in candidiasis is typically thick, white, and clumpy, resembling cottage cheese, with no odor. The vaginal pH remains normal (less than 4.5), and microscopic examination of a potassium hydroxide (KOH) wet mount reveals fungal hyphae, pseudohyphae, or budding yeast cells.
    \item \textbf{Chlamydia (Chlamydia trachomatis):} A bacterial sexually transmitted infection that primarily infects the endocervix in women and the urethra in men. While frequently asymptomatic, it can cause dysuria, post-coital bleeding, and mucopurulent cervicitis. Unlike trichomoniasis, chlamydia does not cause diffuse vaginal inflammation or frothy discharge, and it must be diagnosed using molecular assays.
    \item \textbf{Gonorrhea (Neisseria gonorrhoeae):} A bacterial STI causing acute, purulent urethritis in men and cervicitis in women. In women, the discharge is typically thick, purulent, and yellow-white, arising from the cervical os rather than the vaginal walls. In men, gonorrhea causes severe, acute dysuria and profuse penile discharge, whereas trichomoniasis-induced urethritis is generally milder.
    \item \textbf{Genital Herpes (Herpes Simplex Virus):} Primary or recurrent HSV outbreaks can cause severe vulvar pain, burning, dysuria, and vaginal discharge. However, HSV is distinguished by the presence of painful, vesicular lesions that rupture to form shallow, exquisitely tender ulcers, which are absent in uncomplicated trichomoniasis.
    \item \textbf{Non-Infectious Vaginitis:} This includes contact dermatitis (allergic or irritant reactions to soaps, douches, latex condoms, or spermicides) and atrophic vaginitis. Atrophic vaginitis occurs in postmenopausal women due to estrogen deficiency, leading to thinning of the vaginal epithelium, dryness, dyspareunia, and elevated vaginal pH, which can be misdiagnosed as trichomoniasis.
\end{itemize}

\section*{When to Seek Medical Care}
Individuals should seek medical evaluation from a general practitioner, gynecologist, or sexual health clinic if they experience any signs or symptoms suggestive of a urogenital infection. These symptoms include unusual vaginal discharge (especially if it is frothy, yellow-green, or foul-smelling), abnormal penile discharge, persistent itching, burning or irritation in the genital area, pain or discomfort during urination, or pain during sexual intercourse. Furthermore, anyone whose sexual partner has been diagnosed with trichomoniasis must seek medical evaluation and treatment immediately, even if they have no symptoms themselves, to prevent re-infecting their partner.

While trichomoniasis itself is not a life-threatening condition, it can lead to complications that require urgent or emergency medical care, particularly in pregnant women or when the infection ascends to the upper reproductive tract. Emergency medical services (such as calling local emergency services in many countries, 911 in the United States, or 999 in the United Kingdom) or immediate emergency department evaluation should be sought if a person experiences:
\begin{itemize}
    \item Severe, acute lower abdominal or pelvic pain.
    \item High fever, chills, nausea, and vomiting, which may indicate acute Pelvic Inflammatory Disease (PID) or systemic sepsis.
    \item In pregnant individuals, signs of preterm labor (such as regular uterine contractions or abdominal cramping before 37 weeks of gestation) or premature rupture of membranes (a sudden gush or continuous trickle of fluid from the vagina).
    \item Signs of a severe allergic reaction (anaphylaxis) to prescribed medications, such as hives, swelling of the face, lips, or throat, and difficulty breathing.
\end{itemize}

\section*{Diagnosis and Investigations}
Accurate diagnosis of trichomoniasis cannot be made solely on clinical grounds due to the low specificity of signs and symptoms. Laboratory investigation is required to confirm the presence of \textit{Trichomonas vaginalis}. The diagnostic process involves:
\begin{itemize}
    \item \textbf{Clinical Evaluation:} A detailed sexual history is taken, followed by a physical examination. In females, a speculum examination is performed to visualize the vagina and cervix, checking for characteristic discharge, mucosal inflammation, and the classic ``strawberry cervix.'' In males, the external genitalia are examined for signs of urethral discharge or inflammation.
    \item \textbf{Vaginal pH Measurement:} During the clinical exam, vaginal fluid pH is measured using pH indicator paper. A pH greater than 4.5 is highly suggestive of trichomoniasis or bacterial vaginosis, though not diagnostic on its own.
    \item \textbf{Diagnostic Modalities:}
    \begin{itemize}
        \item \textbf{Wet Mount Microscopy:} This is the traditional and most widely used point-of-care test. A swab of vaginal discharge is mixed with a drop of normal saline on a glass slide and immediately examined under a light microscope. The diagnosis is confirmed by visualizing motile, pear-shaped trichomonads exhibiting a characteristic jerky, twitching motion, along with numerous polymorphonuclear leukocytes. While wet mount is inexpensive and provides immediate results, its sensitivity is poor, ranging from 50\% to 60\%, meaning that a negative wet mount does not rule out infection. Its sensitivity is even lower in males.
        \item \textbf{Nucleic Acid Amplification Tests (NAATs):} NAATs represent the current gold standard for diagnosing trichomoniasis due to their superior sensitivity (often greater than 95\% to 100\%) and specificity. These molecular tests detect the DNA or RNA of \textit{T. vaginalis}. NAATs can be performed on vaginal, cervical, or urethral swabs in women, urethral swabs in men, and first-void urine specimens in both men and women. Because of their high accuracy, NAATs are highly recommended to confirm diagnosis, especially when microscopy is negative but clinical suspicion remains high.
        \item \textbf{Rapid Antigen Tests and Molecular Probes:} These are FDA-approved point-of-care tests that detect specific protozoan antigens or nucleic acids. They offer high sensitivity and specificity (approaching that of laboratory NAATs) and provide results in 10 to 30 minutes, making them valuable in settings without access to immediate microscopy or complex laboratory infrastructure.
        \item \textbf{Protozoal Culture:} Culture involves inoculating vaginal or urethral specimens into specialized liquid media (such as Diamond's medium or the InPouch TV system) and incubating them for 2 to 7 days. The media is examined microscopically every 24 to 48 hours for the growth of trichomonads. While culture was historically the gold standard, it has been largely superseded by NAATs due to the long turnaround time and specialized storage requirements, though it remains useful for evaluating drug-resistant strains.
    \end{itemize}
    \item \textbf{Screening Recommendations:} The Centers for Disease Control and Prevention (CDC) recommends annual screening for trichomoniasis in all HIV-infected women due to the high rate of co-infection and the risk of rapid HIV disease progression. Screening should also be considered for women in high-prevalence settings (such as correctional facilities or STI clinics) or those with multiple sexual partners.
\end{itemize}

\section*{Management}
The primary goal of management is to cure the infection, alleviate symptoms, and prevent transmission or re-infection. Because trichomoniasis is a systemic infection that colonizes tissues beyond the vagina (such as the urethra and periurethral glands), topical treatments (like vaginal creams or suppositories) are ineffective and should not be used. Systemic antibiotic therapy is mandatory.

The standard pharmacological treatments are nitroimidazole compounds:
\begin{itemize}
    \item \textbf{Pharmacotherapy Regimens:}
    \begin{itemize}
        \item \textbf{Metronidazole:} For females, the current gold standard recommended by the CDC is metronidazole 500 mg orally twice daily for 7 days. Clinical trials have demonstrated that this multi-day regimen significantly reduces treatment failure and recurrence compared to the historical single-dose regimen. For males, a single oral dose of metronidazole 2g is still considered an acceptable option, though a 7-day course (500 mg twice daily) may also be prescribed.
        \item \textbf{Tinidazole:} An alternative first-line therapy is tinidazole administered as a single 2g oral dose. Tinidazole has a longer half-life than metronidazole and is generally associated with fewer gastrointestinal side effects, but it is typically more expensive.
    \end{itemize}
    \item \textbf{Adverse Effects and Alcohol Restriction:} Common side effects of nitroimidazoles include nausea, abdominal pain, a metallic taste in the mouth, dizziness, and headache. Crucially, patients must be instructed to strictly avoid alcohol consumption during treatment and for a period afterward (24 hours for metronidazole; 72 hours for tinidazole). Alcohol ingestion during this window can trigger a disulfiram-like reaction, characterized by severe vomiting, abdominal cramps, flushing, tachycardia, and headaches.
    \item \textbf{Management of Sexual Partners:} To prevent reinfection (often referred to as ``ping-pong'' transmission), all current sexual partners of the index patient must be treated simultaneously, regardless of whether they have symptoms or positive test results. Clinicians may utilize Expedited Partner Therapy (EPT), where legally permissible, to provide prescriptions or medications for the patient's partner without a formal clinical evaluation.
    \item \textbf{Abstinence Period:} Patients and their partners must completely abstain from unprotected sexual intercourse (including vaginal, anal, and oral sex) until both have completed their full course of antibiotics and are entirely asymptomatic. This typically requires abstaining for at least 7 days after completing a single-dose regimen, or until the end of a multi-day regimen.
    \item \textbf{Special Populations:}
    \begin{itemize}
        \item \textbf{Pregnancy:} Metronidazole is considered safe to use in all trimesters of pregnancy. The recommended regimen is metronidazole 500 mg twice daily for 7 days. Treatment is critical to mitigate the risk of adverse birth outcomes. Tinidazole should be avoided during pregnancy due to a lack of sufficient safety data.
        \item \textbf{Lactation:} Both metronidazole and tinidazole are excreted in breast milk. To minimize infant exposure, breastfeeding should be withheld for 12 to 24 hours after a maternal dose of metronidazole, and for 72 hours after a dose of tinidazole.
    \end{itemize}
    \item \textbf{Refractory and Resistant Cases:} Treatment failure is most commonly due to non-adherence or re-infection from an untreated partner. However, low-level resistance to metronidazole does occur. If non-adherence and re-infection are ruled out, resistant cases are managed with a high-dose course of metronidazole or tinidazole (e.g., 2g orally daily for 7 days) in consultation with an infectious disease specialist.
\end{itemize}

\section*{Complications}
When left untreated, trichomoniasis can lead to significant physical and psychosocial complications. The chronic inflammatory response induced by the parasite can cause long-term damage to the reproductive tract:
\begin{itemize}
    \item \textbf{Increased Risk of HIV Acquisition and Transmission:} This is one of the most critical public health complications of trichomoniasis. The localized inflammation recruits CD4+ T-lymphocytes and macrophages (the primary target cells for HIV) to the vaginal and cervical mucosa. Concurrently, the parasite causes micro-ulcerations in the epithelium, breaching the physical barrier against viral entry. Consequently, individuals with trichomoniasis are 2 to 3 times more likely to acquire HIV if exposed. Similarly, HIV-positive individuals co-infected with trichomoniasis shed higher levels of HIV in their genital secretions, increasing their infectivity to others.
    \item \textbf{Adverse Pregnancy Outcomes:} Pregnant women with untreated trichomoniasis are at a significantly higher risk for obstetric complications, including:
    \begin{itemize}
        \item \textbf{Preterm Delivery:} Birth occurring before 37 weeks of gestation, which can lead to neonatal respiratory distress and developmental delays.
        \item \textbf{Low Birth Weight:} The infant weighing less than 2500 grams at birth.
        \item \textbf{Premature Rupture of Membranes (PROM):} Rupture of the amniotic sac before the onset of labor.
    \end{itemize}
    \item \textbf{Pelvic Inflammatory Disease (PID):} Although less frequently linked to PID than chlamydia or gonorrhea, \textit{T. vaginalis} can ascend into the uterus and fallopian tubes, causing endometritis and salpingitis. This upper tract inflammation can lead to permanent tubal scarring, increasing the risk of ectopic pregnancy, chronic pelvic pain, and tubal factor infertility.
    \item \textbf{Cervical Cancer Association:} Chronic inflammation and cellular changes in the cervix associated with long-term trichomoniasis are linked to an increased risk of cervical neoplasia and cervical cancer. The infection is also associated with an increased persistence of high-risk Human Papillomavirus (HPV) strains.
    \item \textbf{Complications in Males:} In men, chronic infection can lead to urethral strictures, epididymitis, and prostatitis. Chronic prostatic inflammation caused by \textit{T. vaginalis} has also been investigated as a potential risk factor for the development of prostate cancer and male infertility due to impaired sperm motility and viability.
\end{itemize}

\section*{Prognosis and Outcomes}
The prognosis for individuals diagnosed with trichomoniasis is excellent when the infection is identified promptly and treated with appropriate systemic antibiotic therapy. Cure rates exceed 90\% to 95\% with standard regimens, provided the patient adheres to the medication schedule and avoids alcohol. Following successful treatment, the urogenital mucosa heals, vaginal pH returns to normal, and there are typically no long-term physical sequelae.

However, the rate of re-infection is high, with studies showing that up to 5\% to 20\% of individuals become re-infected within three months of treatment. This high recurrence rate is almost exclusively due to re-exposure from untreated or inadequately treated sexual partners, highlighting the critical importance of partner notification and simultaneous treatment. To address this issue, the CDC and other major health guidelines recommend that all sexually active women treated for trichomoniasis be re-tested (preferably using NAATs) approximately three months after their initial treatment, regardless of whether they believe their partners were treated or if they remain asymptomatic.

If left untreated, the infection can persist for months or even years, acting as a chronic, low-grade source of inflammation and a continuous reservoir for transmission. Chronic untreated cases are more likely to result in the upper genital tract complications and adverse obstetric outcomes detailed in the previous section.

\section*{Prevention and Self-Care}
Preventing the acquisition and transmission of trichomoniasis relies on a combination of barrier protection, regular screening, open partner communication, and strict adherence to self-care protocols during treatment:
\begin{itemize}
    \item \textbf{Barrier Methods:} The most effective way to prevent transmission of trichomoniasis and other STIs during sexual activity is the consistent and correct use of male latex condoms or female polyurethane condoms during vaginal intercourse. Dental dams should be used to prevent transmission during oral-vaginal contact.
    \item \textbf{Limiting Sexual Partners:} Reducing the number of sexual partners and choosing a mutually monogamous relationship with an uninfected partner significantly lowers the risk of exposure.
    \item \textbf{Routine Screening and Testing:} Individuals who are sexually active, particularly those with new or multiple partners, should undergo regular STI screenings. Because trichomoniasis is often asymptomatic, screening is the only way to detect and treat silent infections.
    \item \textbf{Avoiding Vaginal Douching:} Douching should be avoided entirely, as it washes away protective bacteria and disrupts the natural, acidic pH of the vagina, increasing susceptibility to \textit{T. vaginalis} and other infections.
    \item \textbf{Partner Notification and Treatment:} If diagnosed, patients must notify all recent sexual partners so they can be tested and treated simultaneously. This is critical to break the cycle of re-infection.
    \item \textbf{Self-Care During Treatment:}
    \begin{itemize}
        \item \textbf{Complete the Medication:} Patients must take the entire prescribed course of antibiotics, even if their symptoms disappear before the medication is finished. Stopping early can lead to treatment failure and contribute to antibiotic resistance.
        \item \textbf{Abstain from Sexual Activity:} Complete sexual abstinence must be maintained by both the patient and their partner(s) until both have completed the treatment course and all symptoms have resolved.
        \item \textbf{Avoid Alcohol:} Strictly avoid alcohol during metronidazole therapy and for 24 hours after completion, or during tinidazole therapy and for 72 hours after completion, to prevent a dangerous disulfiram-like reaction.
        \item \textbf{Maintain Genital Hygiene:} Keep the genital area clean and dry. Avoid using scented soaps, bubble baths, feminine hygiene sprays, or colored toilet paper, which can irritate the sensitive, inflamed mucosal tissues. Wear loose-fitting, breathable cotton underwear.
    \end{itemize}
\end{itemize}

\section*{Sources and Further Reading}
\begin{itemize}
    \item \textbf{World Health Organization (WHO) - Sexually Transmitted Infections Treatment Guidelines:} Provides global guidelines, epidemiological updates, and public health strategies. URL: \texttt{https://www.who.int/teams/global-hiv-hepatitis-and-stis-programmes/stis/treatment-guidelines}
    \item \textbf{Centers for Disease Control and Prevention (CDC) - Sexually Transmitted Infections Treatment Guidelines:} The primary clinical reference for diagnosis and management in North America. URL: \texttt{https://www.cdc.gov/std/treatment-guidelines/trichomoniasis.htm}
    \item \textbf{British Association for Sexual Health and HIV (BASHH) - National Guideline for the Management of Trichomonas vaginalis:} UK clinical guidelines covering diagnostic and treatment pathways. URL: \texttt{https://www.bashh.org/guidelines}
    \item \textbf{American Sexual Health Association (ASHA) - Patient Resources on Trichomoniasis:} A consumer-focused resource offering clear educational materials on STIs. URL: \texttt{https://www.ashasexualhealth.org/stdsstis/trichomoniasis/}
    \item \textbf{Mayo Clinic - Trichomoniasis Information:} A comprehensive patient-oriented database detailing symptoms, causes, and treatment options. URL: \texttt{https://www.mayoclinic.org/diseases-conditions/trichomoniasis/symptoms-causes/syc-20378443}
    \item \textbf{Australian Sexually Transmitted Infections (STI) Guidelines - Trichomoniasis:} National guidelines for Australian clinicians regarding screening, management, and contact tracing. URL: \texttt{https://www.sti.guidelines.org.au/sexually-transmissible-infections/trichomoniasis}
\end{itemize}